\section*{Chapter \ref{ch:b3:residues} --- Laurent Series and
the Residue Theorem}

\begin{solution}{exo:b3:residues:1}
$\frac{\sin z}z = 1 - \frac{z^2}6 + \cdots$: removable,
residue $0$. $\frac{\eu^z - 1}{z^2} = \frac1z + \frac12 +
\frac z6 + \cdots$: simple pole, residue $1$.
$\frac1{z(z-1)^2}$: simple pole at $0$ with residue
$\frac1{(0-1)^2} = 1$; double pole at $1$ with residue
$\frac{\dd}{\dd z}\bigl(\frac1z\bigr)\big|_{z=1} = -1$.
$\frac{\cos z}{z^3} = \frac1{z^3} - \frac1{2z} + \cdots$: pole
of order $3$, residue $-\frac12$. $\eu^{1/z} = \sum_{n\geq0}
\frac{z^{-n}}{n!}$: essential, residue $1$. $\frac1{\sin z}$:
simple poles at $n\pi$; residue $\frac1{\cos 0} = 1$ at $0$,
$\frac1{\cos\pi} = -1$ at $\pi$
(\cref{met:b3:residues:compute}, $g/h'$).
\end{solution}

\begin{solution}{exo:b3:residues:2}
With $z = \eu^{\iu t}$, $\cos t = \frac{z + z^{-1}}2$, $\dd t
= \frac{\dd z}{\iu z}$:
\[
\int_0^{2\pi}\frac{\dd t}{a + \cos t}
= \oint_{\abs z = 1}\frac{2\,\dd z}{\iu\,(z^2 + 2az + 1)} .
\]
The roots $z_\pm = -a \pm \sqrt{a^2 - 1}$ satisfy $z_+z_- =
1$ with $\abs{z_+} < 1 < \abs{z_-}$; the residue at $z_+$ is
$\frac{1}{z_+ - z_-} = \frac1{2\sqrt{a^2-1}}$, so the integral
is $\frac2\iu\cdot2\iu\pi\cdot\frac1{2\sqrt{a^2-1}} =
\frac{2\pi}{\sqrt{a^2-1}}$. As $a \to 1^+$ it blows up (the
integrand peaks at $t = \pi$); as $a \to \infty$ it behaves
like $\frac{2\pi}a$, matching $\int\frac{\dd t}a$.
\end{solution}

\begin{solution}{exo:b3:residues:3}
First integral: upper poles of $\frac{z^2}{1+z^6}$ at $p =
\eu^{\iu\pi/6}, \iu, \eu^{5\iu\pi/6}$; at each,
$\operatorname{Res} = \frac{p^2}{6p^5} = \frac{p^3}{6p^6} =
-\frac{p^3}6$, and $p^3$ takes the values $\iu, -\iu, \iu$:
sum of residues $-\frac{\iu}{6}$. The arc is
$O(R^{-4})\cdot O(R) \to 0$:
\[
\int_\R\frac{x^2\,\dd x}{1 + x^6} =
2\iu\pi\Bigl(-\frac \iu6\Bigr) = \frac\pi3 .
\]
Second: double pole at $\iu$ of $\frac1{(1+z^2)^2} =
\frac1{(z-\iu)^2(z+\iu)^2}$:
\[
\operatorname{Res} = \frac{\dd}{\dd z}\,(z +
\iu)^{-2}\Big|_{z=\iu} = \frac{-2}{(2\iu)^3} =
\frac{-2}{-8\iu} = -\frac\iu4,
\qquad
\int_\R\frac{\dd x}{(1+x^2)^2} =
2\iu\pi\cdot\Bigl(-\frac\iu4\Bigr) = \frac\pi2 ,
\]
consistent with \cref{exo:b3:fouriertransform:3}(b).
\end{solution}

\begin{solution}{exo:b3:residues:4}
Close $\frac{\eu^{\iu tz}}{z^2 + a^2}$ in the upper half-plane
($t \geq 0$): there $\abs{\eu^{\iu tz}} =
\eu^{-t\operatorname{Im}z} \leq 1$, so the arc contributes
$O(R^{-2})\cdot O(R) \to 0$. The single enclosed pole $\iu a$
is simple with residue $\frac{\eu^{-at}}{2\iu a}$:
\[
\int_\R\frac{\eu^{\iu tx}}{x^2 + a^2}\dd x =
\frac{\pi}{a}\,\eu^{-at},
\qquad\text{hence}\qquad
\int_\R\frac{\cos(tx)}{x^2+a^2}\dd x = \frac\pi
a\,\eu^{-a\abs t}
\]
(real part; even in $t$). This is the inversion counterpart
of $\widehat{\eu^{-a\abs x}} = \frac{2a}{a^2+\xi^2}$
(\cref{exo:b3:fouriertransform:1}): the two computations
confirm each other through
\cref{thm:b3:fouriertransform:inversion}.
\end{solution}

\begin{solution}{exo:b3:residues:5}
Partial fractions: $f = \frac1{z-2} - \frac1{z-1}$.
On $\abs z < 1$ (Taylor):
$f = \sum_{n\geq0}\bigl(1 - 2^{-n-1}\bigr)z^n$.
On $1 < \abs z < 2$: $\frac1{z-2} =
-\sum_{n\geq0}\frac{z^n}{2^{n+1}}$ and $-\frac1{z-1} =
-\sum_{n\geq1}z^{-n}$: a genuine two-sided series.
On $\abs z > 2$: $f = \sum_{n\geq1}\bigl(2^{n-1} -
1\bigr)z^{-n}$. The three differ because Laurent expansions
are attached to \emph{annuli}, not points: each region has its
own geometric expansions. The $z^{-1}$-coefficients ($-1$ and
$0$ respectively) are integrals over circles encircling
\emph{the singularities inside}, not residues at $0$ ($f$ is
holomorphic at $0$): for $1 < \abs z < 2$ the coefficient
$-1$ is $\operatorname{Res}(f, 1)$; for $\abs z > 2$ the
coefficient $0$ is $\operatorname{Res}(f,1) +
\operatorname{Res}(f,2) = -1 + 1$. The residues of $f$: $-1$
at $1$ and $+1$ at $2$.
\end{solution}

\begin{solution}{exo:b3:residues:6}
(a) The Laurent series $\sum_nz^{-n}/n!$ has infinitely many
negative terms: essential. Solving $\eu^{1/z} = w$ ($w \neq
0$): $\frac1z = \log\abs w + \iu\arg w + 2\iu\pi k$, so
\[
z_k = \frac1{\log\abs w + \iu\arg w + 2\iu\pi k}
\xrightarrow[k\to\infty]{} 0 :
\]
every nonzero value is attained infinitely often near $0$ ---
stronger than density.
(b) Along $z = 1/x$, $x \to +\infty$: $\eu^x \to \infty$;
along $z = \iu/y$: modulus $1$. A pole requires
$\abs{f(z)} \to \infty$ \emph{along every approach}: here the
limit simply does not exist, even in $[0, +\infty]$.
\end{solution}

\begin{solution}{exo:b3:residues:7}
(a) On $\abs z = 1$: $\abs{z^7 + z - 1} \leq 3 < 4 =
\abs{-4z^3}$. Rouch\'e with $f = -4z^3$, $g = z^7 + z - 1$:
three zeros in the disc.
(b) On $\abs z = 1$: $\abs{z^4 + 1} \leq 2 < 5 = \abs{5z}$:
one zero in $\abs z < 1$. On $\abs z = 2$: $\abs{5z + 1} \leq
11 < 16 = \abs{z^4}$: four zeros in $\abs z < 2$. Hence three
zeros in the annulus.
(c) For $P = z^n + a_{n-1}z^{n-1} + \dots$: on $\abs z = R >
1 + \sum\abs{a_k}$, $\abs{P - z^n} \leq
\bigl(\sum\abs{a_k}\bigr)R^{n-1} < R^n = \abs{z^n}$: $P$ has
exactly $n$ zeros in $D(0, R)$ --- d'Alembert--Gauss with
multiplicity, by pure counting.
\end{solution}

\begin{solution}{exo:b3:residues:8}
(a) Suppose $f(a) = 0$, $f \not\equiv 0$: choose $r$ with $f$
zero-free on the circle $C = \partial D(a, r)$ (isolated
zeros) and $m = \min_C\abs f > 0$. By
\cref{thm:b3:holomorphic:weierstrassconv}, $f_n \to f$ and
$f_n' \to f'$ uniformly on $C$; for large $n$, $\abs{f_n}
\geq m/2$ on $C$, so
\[
\frac1{2\iu\pi}\int_C\frac{f_n'}{f_n}
\longrightarrow \frac1{2\iu\pi}\int_C\frac{f'}{f} \geq 1
\]
(the limit counts the zero $a$; convergence because
numerators converge uniformly and denominators are uniformly
bounded below). The left side is an integer counting zeros of
$f_n$ in the disc: it must be $\geq 1$ eventually ---
contradicting zero-freeness. So $f$ is zero-free.

(b) Fix $w \in \Omega$ and apply (a) on the connected open set
$\Omega\setminus\{w\}$ (removing a point of an open connected
subset of $\C$ preserves connectedness) to $g_n(z) = f_n(z) -
f_n(w)$, zero-free there by injectivity, converging to $g = f
- f(w)$. If $f$ is nonconstant, $g \not\equiv 0$ on
$\Omega\setminus\{w\}$, so $g$ is zero-free there: $f(z) \neq
f(w)$ for all $z \neq w$. As $w$ was arbitrary, $f$ is
injective.
\end{solution}

\begin{solution}{exo:b3:residues:9}
The sector boundary consists of $[0, R]$, the arc $A_R$, and
the ray $\eu^{2\iu\pi/n}[0, R]$ reversed. Inside lies the
single pole $p = \eu^{\iu\pi/n}$ of $\frac1{1+z^n}$, with
residue $\frac1{np^{n-1}} = \frac{p}{np^n} = -\frac pn$. On
the return ray, $z = \eu^{2\iu\pi/n}x$ gives $z^n = x^n$ and
$\dd z = \eu^{2\iu\pi/n}\dd x$; the arc is
$O(R^{-n})\cdot O(R) \to 0$. Hence
\[
\bigl(1 - \eu^{2\iu\pi/n}\bigr)
\int_0^\infty\frac{\dd x}{1 + x^n}
= 2\iu\pi\Bigl(-\frac{\eu^{\iu\pi/n}}n\Bigr),
\quad\text{so}\quad
\int_0^\infty\frac{\dd x}{1+x^n}
= \frac{2\iu\pi}{n\,\bigl(\eu^{\iu\pi/n} -
\eu^{-\iu\pi/n}\bigr)} = \frac{\pi}{n\sin(\pi/n)} .
\]
$n = 2$: $\frac\pi{2\sin(\pi/2)} = \frac\pi2 =
[\arctan]_0^\infty$. As $n \to \infty$: the value tends to
$1$, and indeed the integrand tends to $\mathbf 1_{\intco01}$
(with DCT domination $\min(1, x^{-2})$ for $n \geq 2$).
\end{solution}

\begin{solution}{exo:b3:residues:10}
On $\abs z = R$ large, $\abs f \leq C R^{-2}$: $\abs{\oint_{C_R}
f} \leq 2\pi R\cdot CR^{-2} \to 0$. But for $R$ beyond all
poles, the residue theorem gives $\oint_{C_R}f =
2\iu\pi\sum_{\text{all }p}\operatorname{Res}(f, p)$: the total
sum vanishes. For $f = \frac1{z(z-1)(z-2)}$: residues
$\frac1{(-1)(-2)} = \frac12$ at $0$, $\frac1{1\cdot(-1)} = -1$
at $1$, $\frac1{2\cdot1} = \frac12$ at $2$ --- summing to $0$
as predicted, and
\[
\frac1{z(z-1)(z-2)} = \frac{1/2}{z} - \frac1{z - 1} +
\frac{1/2}{z-2} :
\]
the residues \emph{are} the partial fraction coefficients, and
the zero-sum identity supplies a free consistency check (or
determines the last coefficient from the others).
\end{solution}

\begin{solution}{pb:b3:residues:1}
\textbf{1.} $\sin(\pi z)$ has simple zeros exactly at $\Z$
($\sin\pi z = 0$ iff $z \in \Z$, and $(\sin\pi z)' = \pi\cos\pi
z \neq 0$ there), and $\cos(\pi n) \neq 0$: $\pi\cot(\pi z) =
\pi\cos(\pi z)/\sin(\pi z)$ has simple poles at $\Z$ with
\[
\operatorname{Res}(\pi\cot\pi z,\ n)
= \frac{\pi\cos(\pi n)}{\pi\cos(\pi n)} = 1
\]
($g/h'$ rule, \cref{met:b3:residues:compute}).

\textbf{2.} $\frac{\sin(\pi z)}{\pi z} = 1 - \frac{(\pi z)^2}6
+ \frac{(\pi z)^4}{120} - \cdots$ is holomorphic and nonzero
near $0$: its reciprocal is holomorphic
(\cref{def:b3:holomorphic:holo}: quotient), with series $1 +
\frac{(\pi z)^2}{6} + \frac{7(\pi z)^4}{360} + \cdots$
(identify: $(1 - u)^{-1}$-style coefficients from $u =
\frac{(\pi z)^2}6 - \frac{(\pi z)^4}{120}$: the $z^4$
coefficient is $\frac1{36} - \frac1{120} = \frac{7}{360}$).
Multiply by $\cos(\pi z) = 1 - \frac{(\pi z)^2}2 + \frac{(\pi
z)^4}{24} - \cdots$ and divide by $z$:
\[
\pi\cot(\pi z) = \frac1z\Bigl[1 + \pi^2z^2\Bigl(\frac16 -
\frac12\Bigr) + \pi^4z^4\Bigl(\frac7{360} - \frac1{12} +
\frac1{24}\Bigr)\Bigr] + \cdots
= \frac1z - \frac{\pi^2}3\,z - \frac{\pi^4}{45}\,z^3 -
\cdots
\]
($\frac7{360} - \frac{30}{360} + \frac{15}{360} =
-\frac8{360} = -\frac1{45}$).

\textbf{3.} Vertical sides $z = \pm(N + \frac12) + \iu y$:
by $\pi$-periodicity of $\cot$, $\cot(\pi z) = \cot(\pm\frac\pi2
+ \iu\pi y) = -\tan(\iu\pi y) = -\iu\tanh(\pi y)$, of modulus
$\leq 1$. Horizontal sides $z = x \pm \iu(N + \frac12)$: from
$\abs{\cot(a + \iu b)}^2 = \frac{\cos^2a +
\sinh^2b}{\sin^2a + \sinh^2b} \leq \frac{1 +
\sinh^2b}{\sinh^2b} = \coth^2 b$,
\[
\abs{\cot(\pi z)} \leq \coth\bigl(\pi(N + \tfrac12)\bigr)
\leq \coth(\pi/2) < 1.1 .
\]
Both bounds are $\leq 2$.

\textbf{4.} Apply \cref{thm:b3:residues:residue} to $F(z) =
\pi\cot(\pi z)f(z)$ on $C_N$ ($N$ beyond all poles of $f$):
\[
\frac1{2\iu\pi}\oint_{C_N}F = \sum_{n=-N}^{N}f(n) +
\sum_p\operatorname{Res}(F, p),
\]
the integer poles contributing $f(n)$ (question 1; $f$
holomorphic there). On $C_N$: $\abs F \leq 2\pi\cdot
C\abs z^{-2} \leq 2\pi C N^{-2}$, and the perimeter is $8(N +
\frac12)$: the integral is $O(1/N) \to 0$. Let $N \to
\infty$: the displayed summation formula.

\textbf{5.} The decay $\abs f = O(\abs z^{-2})$ killed the
contour integral \emph{and} made $\sum\abs{f(n)}$ converge.
For $f(z) = \frac1{z + \frac12}$: the symmetric sums
$\sum_{-N}^N\frac1{n + \frac12}$ telescope to $0$ (the terms
$n$ and $-n - 1$ cancel), and the right side is
$-\operatorname{Res}\bigl(\frac{\pi\cot\pi z}{z + \frac12},
-\frac12\bigr) = -\pi\cot(-\frac\pi2) = 0$: consistent ---
but $\sum\abs{f(n)}$ diverges; the method computes the
symmetric (principal value) limit only.

\textbf{6.} $g(z) = \frac{\pi\cot(\pi z)}{z^2}$: poles at the
nonzero integers with residues $\frac1{n^2}$, and at $0$
where, by question 2,
\[
g(z) = \frac1{z^3} - \frac{\pi^2}{3z} - \frac{\pi^4}{45}z -
\cdots :
\qquad \operatorname{Res}(g, 0) = -\frac{\pi^2}3 .
\]
The contour integral over $C_N$ tends to $0$ as in question 4
($\abs{g} = O(N^{-2})$ on $C_N$). Hence $0 =
\sum_{n\neq0}\frac1{n^2} - \frac{\pi^2}3$: $2\zeta(2) =
\frac{\pi^2}3$, $\zeta(2) = \frac{\pi^2}6$.

\textbf{7.} $g(z) = \frac{\pi\cot(\pi z)}{z^4} = \frac1{z^5}
- \frac{\pi^2}{3z^3} - \frac{\pi^4}{45z} - \cdots$: residue
at $0$ equal to $-\frac{\pi^4}{45}$, and $0 =
2\zeta(4) - \frac{\pi^4}{45}$: $\zeta(4) =
\frac{\pi^4}{90}$.

\textbf{8.} With $g = \pi\cot(\pi z)/z^{2k}$: the residue at
$0$ is the coefficient $a_{2k-1}$ of $z^{2k-1}$ in the
expansion of $\pi\cot(\pi z)$, and the vanishing contour gives
$2\zeta(2k) + a_{2k-1} = 0$: $\zeta(2k) = -a_{2k-1}/2$, a
rational multiple of $\pi^{2k}$ since the cotangent's
coefficients are. One more division step yields $a_5 =
-\frac{2\pi^6}{945}$, whence $\zeta(6) = \frac{\pi^6}{945}$.
For $\zeta(3)$: the natural kernel $g = \pi\cot(\pi z)/z^3$
produces $\sum_{n\neq0}\frac1{n^3} = 0$ by oddness --- the
method proves $0 = 0$ and is structurally blind to odd zeta
values (no closed form for $\zeta(3)$ is known; its
irrationality, Ap\'ery 1978, needed entirely different
ideas).

\textbf{9.} $F(z) = \frac{\pi\cot(\pi z)}{(z - w)(z + w)}$
has poles at the integers (residues $\frac1{n^2 - w^2}$,
noting the sign: $f(n) = \frac1{(n-w)(n+w)} = \frac1{n^2 -
w^2}$) and simple poles at $\pm w$ with residues
$\frac{\pi\cot(\pm\pi w)}{\pm2w} = \frac{\pi\cot(\pi
w)}{2w}$ each ($\cot$ is odd). Question 4's argument
($\abs f = O(\abs z^{-2})$) gives
\[
\sum_{n\in\Z}\frac1{n^2 - w^2} + \frac{\pi\cot(\pi w)}{w} =
0,
\qquad\text{i.e.}\qquad
\pi\cot(\pi w) = \frac1w + \sum_{n\geq1}\frac{2w}{w^2 -
n^2},
\]
(the $n = 0$ term is $-\frac1{w^2}$; regroup $\pm n$). Normal
convergence on compacts of $\C\setminus\Z$: for $\abs w \leq
R$ and $n \geq 2R$, $\abs{\frac{2w}{w^2 - n^2}} \leq
\frac{2R}{n^2 - R^2} \leq \frac{8R}{3n^2}$.

\textbf{10.} For $\abs w \leq r < 1$: $\frac{2w}{w^2 - n^2} =
-\frac{2w}{n^2}\cdot\frac1{1 - w^2/n^2} =
-2\sum_{k\geq0}\frac{w^{2k+1}}{n^{2k+2}}$, with
$\abs{\text{terms}} \leq 2r^{2k+1}/n^{2k+2}$, summable over
$(n, k)$: Fubini for series rearranges
\[
\pi\cot(\pi w) = \frac1w -
2\sum_{k\geq0}\zeta(2k+2)\,w^{2k+1} .
\]
Matching with question 2: $-2\zeta(2) = -\frac{\pi^2}3$ and
$-2\zeta(4) = -\frac{\pi^4}{45}$ --- the same values. Three
roads to $\frac{\pi^2}6$: Parseval (Fourier series), the
trace of the string's Green operator, and the cotangent's two
expansions; that a sum over frequencies, an operator trace,
and a contour integral agree is no accident --- each is a
face of the same spectral identity.
\end{solution}
